%=========================================================================
%  FEX1001 -- Física Experimental I -- UDESC/CCT
%  Modelo de relatório de laboratório.
%
%  Este arquivo é preenchido automaticamente pela página do laboratório.
%  Os comandos marcados abaixo são substituídos na geração; o restante é
%  seu para editar livremente.
%
%  Compila com pdflatex em qualquer distribuição TeX (inclusive Overleaf):
%  usa apenas a classe article e pacotes padrão.
%=========================================================================
\documentclass[a4paper,11pt]{article}

\usepackage[utf8]{inputenc}
\usepackage[T1]{fontenc}
\usepackage[brazil]{babel}
\usepackage[margin=2.5cm]{geometry}
\usepackage{amsmath}
\usepackage{amssymb}
\usepackage{graphicx}
\usepackage{booktabs}
\usepackage{siunitx}
\usepackage{float}
\usepackage{caption}
\usepackage[hidelinks]{hyperref}

\sisetup{output-decimal-marker={,}, group-separator={.}, group-digits=integer}
% group-digits=integer é essencial: sem isso o siunitx agrupa também as
% casas decimais e 9.79061 sai impresso como "9,790.61".
\captionsetup{font=small,labelfont=bf}

%-------------------------------------------------------------------------
% DADOS DO RELATÓRIO -- preenchidos pela página do laboratório
%-------------------------------------------------------------------------
\newcommand{\experimento}{Experiência}
\newcommand{\subtitulo}{Física Experimental I}
\newcommand{\equipe}{Bancada}
\newcommand{\integrantes}{Nomes da equipe}
\newcommand{\datarelatorio}{\today}

\newcommand{\disciplina}{FEX1001 -- Física Experimental I}
\newcommand{\professor}{Prof. Dr. Rafael C. R. de Lima}
\newcommand{\instituicao}{Universidade do Estado de Santa Catarina --
Centro de Ciências Tecnológicas, Joinville/SC}

%-------------------------------------------------------------------------
\begin{document}

\begin{center}
  {\large\scshape \instituicao}\\[2pt]
  {\small \disciplina \quad|\quad \professor}\\[14pt]
  \rule{\linewidth}{0.5pt}\\[10pt]
  {\LARGE\bfseries \experimento}\\[4pt]
  {\large \subtitulo}\\[10pt]
  \rule{\linewidth}{0.5pt}\\[10pt]
  {\normalsize Equipe (bancada): \equipe}\\[2pt]
  {\normalsize \integrantes}\\[2pt]
  {\small \datarelatorio}
\end{center}

\vspace{6pt}

%=========================================================================
\section{Objetivo}
\input{secoes/objetivo}

%=========================================================================
\section{Observações}
\input{secoes/observacoes}

%=========================================================================
\section{Experimentos}
\input{secoes/experimentos}

\input{tabelas/tabela1_medidas}

%=========================================================================
\section{Teoria}
\input{secoes/teoria}

%=========================================================================
\section{Identificação das variáveis}
\input{tabelas/tabela2_variaveis}

%=========================================================================
\section{Linearização}
\input{secoes/linearizacao}

%=========================================================================
\section{Análise estatística}
\input{tabelas/tabela3_estatistica}

%=========================================================================
\section{Gráfico linear}
\input{tabelas/tabela4_linearizada}

\input{secoes/grafico}

%=========================================================================
\section{Resultados}
\input{secoes/resultados}

%=========================================================================
\section{Conclusões}
\input{secoes/conclusoes}

%=========================================================================
\begin{thebibliography}{9}
\bibitem{halliday} HALLIDAY, D.; RESNICK, R.; WALKER, J.
  \emph{Fundamentos de Física} -- Volume 1: Mecânica.
  8.~ed. Rio de Janeiro: LTC.
\bibitem{gravidade} Aceleração da gravidade local em Joinville/SC.
  Disponível em: \url{https://www.wolframalpha.com}.
\end{thebibliography}

\end{document}
